\section{Two general equality lemmas}
\begin{lemma}\label{intransitive}
Let $H\leq S_k$ be intransitive and generated irredundantly by $k-2$
elements. Then $H=S_A\times S_B$ on exactly two nonempty orbits
$A,B$ partitioning the $k$ letters.
\end{lemma}
\begin{proof}
Successive orbit projections and the normal-quotient normalization
bound give independent length at most $\sum_O(|O|-1)=k-q$, where
$q$ is the number of orbits. Hence $q=2$. Write $a=|A|$, $b=|B|$.
Choose a minimal subfamily $T$ generating the first projection. It
has at most $a-1$ members. Normalize all other generators into the
second-coordinate kernel by words in $T$. They remain independent,
so there are at most $b-1$ of them. Equality forces both bounds to
be attained. Whiston's equality clause makes the first image $S_A$
and the subgroup generated by the normalized extra elements $S_B$.
Thus the whole second factor is in $H$, and the full first image
then gives $H=S_A\times S_B$. An orbit of size one has the trivial
symmetric factor and contributes zero throughout.
\end{proof}

We use the following extremal consequence of UB.Extremal: an independent
string marking of $S_b$ with $b-1$ involutions, for $b\geq7$, is a Coxeter
path. Indeed a sign-modified tree marking cannot consist entirely of
involutions: an edge adjacent to the distinguished edge gives an
element of order three. For the ordinary tree, the string condition
makes its line graph a path, so the tree itself is a path. For $b=5,6$
the complete small certificate verifies the same abstract Coxeter-path
conclusion, including the outer-automorphism image in $S_6$. It also
verifies the trivial centralizer of the subgroup omitting an endpoint.

\begin{lemma}\label{smallfactor}
For $a=2$ or $3$, and $b\geq6$, an involutory string C-marking of
$S_a\times S_b$ of rank $a+b-2$ has two nonempty diagram components.
One generates $S_a$, and the other is the graph of a homomorphism
$S_b\to S_a$; this homomorphism is trivial when $a=3$.
\end{lemma}
\begin{proof}
Choose a minimal subfamily $L$ generating the $S_b$ quotient.
Normal-quotient normalization bounds the remaining family by $a-1$,
so $|L|=b-1$. The images of $L$ are an extremal Coxeter path.
Its noncommuting adjacent vertices must be adjacent in the full
string, so $L$ occupies an interval. Every generator outside that
interval centralizes all but possibly one endpoint of this path.
The endpoint-omission subgroup has trivial centralizer in $S_b$;
therefore every outside generator has trivial $S_b$ image and lies
in the small factor. There are exactly $a-1$ of them.

Put $N=\gen{L}\cap S_a$. Their irredundance relative to $\gen{L}$
implies irredundance relative to $N$ under conjugation by $\gen{L}$.
If $a=2$, the one extra generator forces $N=1$ and is central;
its diagram vertex is isolated, and $\gen L$ is a graph over $S_b$.
If $a=3$, any nontrivial subgroup
of $S_3$, together with a family generating $S_3$, needs at most one
member of that family. Thus two independent extras force $N=1$.
The complement $\gen{L}$ is the graph of a homomorphism
$S_b\to S_3$, whose image is trivial or order two, since $A_b$ is
simple and $b\geq5$. If the image has order two, one of the two
extra involutions lies outside it. That involution and its conjugate
by the image generate $S_3$, making the other extra redundant.
Hence the image is trivial. Now the two small-factor involutions
commute with all of $L$, so the diagram is disconnected.
\end{proof}

\section{The uniform two-large-block obstruction}
Let a transitive string C-group $G$ preserve two blocks of size $k$.
Choose a distinguished generator $r$ interchanging the blocks,
nearest to an endpoint of the string among all such generators.
Identify the blocks through $r$, so $r$ is the pure block swap in
$S_k\wr C_2$. Let $C$ be the other distinguished generators commuting
with $r$, and $R$ the remainder. Only its string neighbors can be in
$R$, so $|R|\leq2$.

The centralizer of $r$ in this wreath product is
$\operatorname{diag}(S_k)\times\gen r$. Multiplying each $c\in C$
by $r$ when necessary normalizes it into the diagonal copy of $S_k$.
The projection is injective on $\gen C$, because the intersection
property gives $r\notin\gen C$. Its image $H$ therefore has the
same involutory string C-marking. Moreover
\[
 \gen C=\{(h,h)r^{\chi(h)}:h\in H\}
\]
for a character $\chi:H\to C_2$. Whiston gives $|C|\leq k-1$,
and consequently $d\leq k+2$.

\begin{theorem}\label{twoblocks}
For $k\geq9$, suppose every transitive string C-group of faithful degree
$k$, other than $S_k,A_k$, has rank at most $k/2+1$. Then a transitive
two-block string C-group on $2k$ letters has $d\leq k$.
\end{theorem}
\begin{proof}
If $d=k+2$, then $|C|=k-1$ and $H=S_k$. Its diagram is a Coxeter
path, so $C$ lies on one side of $r$ and its neighbors: $r$ must be
one place from a string endpoint. The outside neighbor centralizes
all of $C$. But $\ker\chi$ contains $A_k$, whose centralizer in
$S_k$ is trivial. An element $(p,q)r^\epsilon$ centralizing $C$
therefore has $p=q=1$, and is in $\gen r$, a contradiction.

Suppose $d=k+1$. If $r$ is at an endpoint, $|C|=k-1$ and the one
remaining generator centralizes a Coxeter endpoint-omission subgroup
of $C$. Its character kernel contains an $A_{k-1}$ fixing one letter;
that subgroup again has trivial centralizer in $S_k$. The remaining
generator is in $\gen r$, a contradiction. If $r$ is interior but
$|R|=1$, then again $|C|=k-1$ and its diagram is connected. Therefore
$r$ is one place from an endpoint and the sole noncommuting neighbor
is on that outer side. It commutes with all of $C$, so the preceding
$A_k$ centralizer argument excludes it. Thus $|R|=2$ and $|C|=k-2$.

First let $r$ have distance one from an endpoint, so $C$ occupies
one interval (its noncommuting diagram need not be connected).
Its outside neighbor centralizes all of $C$ and does not commute with
$r$. If $H$ is transitive, the hypothesis at degree $k$ and
$k-2>k/2+1$ leave only $S_k,A_k$. This use is legitimate because
the projection onto $H$ preserves the full marking and IP. The same
character-kernel centralizer argument excludes both.
Otherwise Lemma~\ref{intransitive} gives $H=S_A\times S_B$.
If both orbits have size at least four, $\ker\chi$ contains
$A_A\times A_B$ and has trivial centralizer. A singleton orbit is
also excluded by this argument on its complement. The only remaining
possibilities have $a=|A|\in\{2,3\}$ and $b=|B|\geq6$.
If $\chi$ is nontrivial on the large factor, a centralizing
$(p,q)r^\epsilon$ first has $p,q$ supported on $A$ (by $A_B$),
then has $p=q$ by an odd permutation supported on $B$. It commutes
with $r$, a contradiction. If $\chi$ is trivial there,
Lemma~\ref{smallfactor} gives two components. The component at the far
endpoint has extremal Coxeter generators, all odd in its symmetric
factor. A nontrivial restriction of $\chi$ to that component would
make its endpoint a block swap closer to an endpoint than $r$.
Hence $\chi$ is trivial on that far component. Both neighbors of $r$
centralize it: the inner neighbor commutes with all of $C$ except
possibly its first vertex, in the other component. The far component
has $A$ or $B$ as its distinguished orbit; any other nontrivial orbit
has a different size. Its centralizer in $S_k$ therefore preserves
$A,B$. Thus the coordinate permutations of both neighbors preserve
$A,B$. Together with $C$ and $r$ they preserve $A\times\{0,1\}$,
contradicting transitivity.

Finally let $r$ have distance at least two from both endpoints.
Then $C$ occupies two nonempty intervals separated by the missing
neighbors and $r$. Their generated images commute and form a direct
product decomposition of $H$, by the intersection property; we do not
yet assume that either interval has connected noncommuting diagram.
The same smaller-degree hypothesis excludes transitive $H$ other than
$S_k,A_k$, and these cannot be a nontrivial direct product. Thus again
$H=S_A\times S_B$ by Lemma~\ref{intransitive}.
Neither orbit is a singleton. For orbit sizes at least three, the
symmetric factors are centerless and directly indecomposable, so
the two interval factors are precisely those symmetric factors. Their
noncommuting diagrams are therefore connected. If one orbit has
size two, one component is this central $C_2$, and the other is a
graph of a character $S_B\to C_2$. This graph is directly
indecomposable too. Thus in both cases the intervals are exactly the
two noncommuting components, and each attains the maximum independent
length in its symmetric factor.

If $\chi$ is nontrivial on such a component of size at least five,
its extremal Coxeter generators are all odd. Hence the endpoint of
that component at the end of the whole string interchanges the two
large blocks, contradicting the choice of $r$. The same holds for
factors of size two or three. For $S_4$, at least one of the first
two and one of the last two generators is odd: otherwise two even
involutions lie in its abelian normal four-group, and the string
relations make an endpoint generator central in $S_4$.
This again gives a closer block-swapping generator. Therefore
$\chi$ is trivial on both components, and all of $C$ fixes both
large blocks individually.

Each member of $R$ centralizes the opposite component of $C$.
Its coordinate permutations consequently preserve the two sets
$A,B$. This remains true for the possible character graph over the
large symmetric factor when $|A|=2$: its two nontrivial orbits have
different sizes. All of $C$, $R$ and $r$ therefore preserve
$A\times\{0,1\}$ and $B\times\{0,1\}$, contradicting transitivity.
\end{proof}

\section{The finite two-block boundary}
The same argument excludes $d\geq k+2$ for $5\leq k\leq8$.
For $k=5,6$ the complete independent-string certificate proves both
that the extremal markings are Coxeter paths and that the derived
subgroup of an endpoint-omission subgroup has trivial centralizer.
In the outer $S_6$ marking that subgroup is a transitive degree-six
$A_5$, not a point-fixing alternating group. These actual centralizer
checks replace the natural point-fixing description in the endpoint
argument. The cases with $d=k+1$ and $|R|\leq1$ are excluded in the
same way.

For $d=k+1$ and $|R|=2$, the faithful normalized $H\leq S_k$ has
rank $k-2$. If it is intransitive, Lemma~\ref{intransitive} gives
$H=S_a\times S_b$. A singleton orbit is excluded by the alternating
centralizer argument, or by direct indecomposability when the two
intervals are nonempty. The remaining unordered orbit sizes are
$(2,3),(2,4),(3,3),(2,5),(3,4),(2,6),(3,5),(4,4)$.
For a transitive $H$, the accepted primitive bounds and the complete
small transitive-imprimitive census leave only $S_k,A_k$, the
primitive degree-six $S_5$, and the degree-six actions numbered
9 and 11. The first two are excluded as above. In the degree-six
$S_5$ case, its derived subgroup is primitive $A_5$ with trivial
centralizer, so the same argument applies. The degree-eight census
has no rank-six marking; degree five and seven have no nontrivial
imprimitive actions.

For the eight product actions and the two degree-six actions, the
separate, completely proved finite-neighbor lemma enumerates every
$C$-marking, character, compatible split and pair of missing neighbors.
Its 59 normalized markings yield 154 transitive completions, all
redundant. This exact required finite premise therefore excludes all
remaining cases. Consequently the two-large-block bound $d\leq k$
holds for every $k\geq5$. The block sizes three and four are included
in the separately complete degree-six and degree-eight census; they
contain genuine exceptional markings.
